\chapter{Durchführung und Ergebnisse}
\label{chap:kern}
\section{El Salvador -- Bitcoin als Zahlungsmittel}

\subsection{Motivation -- Warum Bitcoin?}

El Salvador wollte mit Bitcoin zwei Probleme lösen: fehlenden Bankzugang für rund zwei Drittel der Bevölkerung und teure Rücküberweisungen aus dem Ausland, die etwa 20\,\% des BIP ausmachten. Präsident Bukele versprach finanzielle Inklusion per Smartphone und günstigere Überweisungen ohne Banken \cite{Alvarez2023}.

\subsection{Umsetzung}

Im September 2021 wurde Bitcoin neben dem US-Dollar gesetzliches Zahlungsmittel -- weltweit erstmalig. Unternehmen mussten Bitcoin akzeptieren; zur Förderung der Nutzung führte die Regierung die Chivo Wallet mit gebührenfreien Transaktionen und einem Startbonus von 30~USD ein. Zudem richtete sie einen Konversionsfonds von rund 150 Millionen USD ein \cite{Alvarez2023}.

\subsection{Ergebnisse -- Die Bitcoin-Einführung in El Salvador}

Für El Salvador wurde Alvarez et al.\ (2023) ausgewertet. Die Studie kombiniert Haushaltsbefragungen, Unternehmensdaten und Blockchain-Daten und zeigt: Das Experiment förderte weder finanzielle Inklusion noch günstigere Überweisungen.

Zwar luden ca.\ 53\,\% der Erwachsenen Chivo herunter, doch nur 20\,\% nutzten die App nach dem Startbonus weiter~\cite{Alvarez2023}. Auch die Neuanmeldungen brachen nach September 2021 ein und lagen ab Anfang 2022 fast bei null (siehe Abb.~\ref{fig:neuanmeldungen}). Die Nutzung hing damit stark vom staatlichen Anreiz ab.

Bitcoin setzte sich auch im Geschäftsverkehr nicht durch: Nur 11,4\,\% der Unternehmen verzeichneten Bitcoin-Verkäufe; 88\,\% davon tauschten die Einnahmen sofort in Dollar um \cite{Alvarez2023} (siehe Abb.~\ref{fig:unternehmen}).

Das Transaktionsvolumen stieg nur zum Start durch den Bonus und brach danach dauerhaft ein (siehe Abb.~\ref{fig:transaktionen}).

Auch das Kernversprechen günstigerer Überweisungen blieb unerfüllt: Im Februar 2022 liefen nur 1,6\,\% der Rücküberweisungen über digitale Wallets~\cite{Alvarez2023}. Gründe waren die starke Bitcoin-Volatilität, technische Probleme der Chivo-App und fehlende Transparenz im Umgang mit öffentlichen Mitteln~\cite{Alvarez2023}.

\section{Vereinigte Arabische Emirate}
 
\subsection{Motivation -- Warum Krypto?}
 
Die VAE wollen ihre Wirtschaft langfristig weniger abhängig vom Öl machen. Seit 2016 setzen sie auf Blockchain und digitale Vermögenswerte, um neue Unternehmen und Investoren anzuziehen~\cite{AlTawil2023}. Die Dubai Blockchain Strategy 2020 verfolgte drei Ziele: effizientere Verwaltung, neue Wirtschaftsbereiche und internationale Sichtbarkeit~\cite{DubaiBlockchain2020}. Das Ziel war dabei nie, eine Kryptowährung als alltägliches Zahlungsmittel einzuführen -- es ging um Standortpolitik und wirtschaftliche Diversifizierung.
 
\subsection{Umsetzung -- Regulierung und Rahmenbedingungen}
 
Die VAE schufen einen regulierten Markt für Kryptowährungen, führten sie aber nicht als Geld ein. 2022 errichtete Dubai mit VARA eine eigene Aufsichtsbehörde für digitale Vermögenswerte~\cite{VARA2025, ElGheriani2025}. 2023 veröffentlichte VARA ein umfassendes Regelwerk, das Lizenzpflichten und Anforderungen an Krypto-Dienstleister festlegt~\cite{VARA2025}. Krypto-Anbieter benötigen eine gültige Lizenz und unterliegen strengen Regeln zur Geldwäschebekämpfung. Anonymitätserhöhende Kryptowährungen sind verboten \cite{VARA2025}. Auch auf Bundesebene fallen Krypto-Anbieter seit 2021 unter das nationale Geldwäschegesetz \cite{AlTawil2023}.
 
Kryptowährungen sind in den VAE kein gesetzliches Zahlungsmittel. Die Zentralbank der VAE stellte ausdrücklich klar, dass sie Krypto-Assets weder als gesetzliches Zahlungsmittel anerkennt noch als Zahlungsmittel akzeptiert -- der Dirham bleibt die einzige offizielle Währung~\cite{CBUAE2025}. Kryptowährungen werden von der Zentralbank ausschließlich als Anlageobjekte mit potenziell hohem Risiko eingestuft~\cite{CBUAE2025}.
 
\subsection{Ergebnisse -- Krypto als Anlage, nicht als Geld}
 
Der Fall der VAE zeigt, dass selbst ein kryptofreundlicher Staat private Kryptowährungen nicht als vollständigen Ersatz für die nationale Währung betrachtet. Die Blockchain-Strategie des Landes führte zu einzelnen praktischen Anwendungen, etwa in den Bereichen Transfers, Immobilien und Tourismus. Auch die Gründung von VARA stärkte die spezialisierte Aufsicht über digitale Vermögenswerte~\cite{DubaiBlockchain2020}~\cite{ElGheriani2025}.

Im alltäglichen Zahlungsverkehr dominiert jedoch weiterhin der Dirham. Kryptozahlungen kommen nur in einzelnen Sektoren vor und werden in der Regel direkt in die nationale Währung umgerechnet. Statt auf private Kryptowährungen setzt die Zentralbank mit dem Digital Dirham auf staatlich kontrolliertes digitales Geld~\cite{CBUAE2025}.

Damit bleiben Kryptowährungen in den VAE vor allem ein Anlage- und Wirtschaftsfaktor, aber kein gewöhnliches Zahlungsmittel.
\section{Europäische Union}

Die Europäische Union (EU) erkennt private Krypto-Assets nicht als gesetzliches Zahlungsmittel an. Stattdessen werden sie auf Grundlage der MiCA-Verordnung als regulierte Vermögenswerte behandelt~\cite{MiCA2023}. Der EU-Fall steht damit für eine Strategie der rechtlichen Einordnung statt monetärer Anerkennung.

\subsection{Motivation -- Schutz des Geldsystems}

Die EU folgt gegenüber Kryptowährungen einer Schutzlogik. In der Eurozone liegt die Geldpolitik bei der Europäischen Zentralbank (EZB). Eine breite Nutzung privater Kryptowährungen als Zahlungsmittel könnte diese Steuerungsfunktion, die Finanzstabilität und die staatliche Kontrolle über Zahlungsströme schwächen.

Deshalb setzt die EU nicht auf eine monetäre Anerkennung, sondern auf rechtliche Einordnung. MiCA schafft einen Rechtsrahmen für Krypto-Assets und Krypto-Dienstleister und legt Anforderungen an Transparenz, Aufsicht, Verbraucherschutz und Marktintegrität fest~\cite{MiCA2023}. Kryptowährungen werden damit als Vermögenswerte erfasst, jedoch nicht als Ersatz für den Euro behandelt. Diese Zurückhaltung zeigt sich auch in der kritischen Bewertung von Bitcoin durch die EZB: Aufgrund begrenzter Zahlungsnutzung, starker Kursschwankungen und technischer Grenzen eignet sich Bitcoin nur eingeschränkt als stabiles Zahlungsmittel~\cite{Bindseil2024}.

\subsection{Umsetzung -- Regulierung statt Anerkennung}

Der europäische Ansatz wird vor allem durch Regulierung umgesetzt, nicht durch die Einführung von Kryptowährungen als Geld. MiCA bindet Krypto-Unternehmen stärker in das europäische Finanzsystem ein, verleiht Krypto-Assets jedoch keinen Status als gesetzliches Zahlungsmittel~\cite{BundesbankMiCAR}.

Parallel arbeitet die EZB am digitalen Euro. Dieser wäre keine private Kryptowährung, sondern Zentralbankgeld in digitaler Form und soll Bargeld ergänzen~\cite{ECB2025}. Dadurch wird die Grenze des europäischen Ansatzes deutlich: Private Krypto-Assets werden reguliert, während neue Geldformen nur unter staatlicher Kontrolle vorbereitet werden.

\subsection{Ergebnisse -- Besitz statt Zahlungsnutzung}

Die empirischen Daten im Euroraum zeigen eine Lücke zwischen Besitz und Zahlungsnutzung. Zamora-Pérez untersucht auf Basis der \emph{Study on the Payment Attitudes of Consumers in the Euro Area} (SPACE Survey) aus dem Jahr 2022 das Verhalten von 39.507 Erwachsenen in 17 Ländern des Euroraums. Die Daten zeigen: Jüngere Gruppen besitzen häufiger Krypto-Assets, aber die Nutzung als Zahlungsmittel bleibt in allen Geburtsgruppen niedriger als der Besitz~\cite{Zamora2026}.

Abbildung~\ref{fig:krypto-besitz-zahlung-eu} zeigt, dass Krypto-Assets im Euroraum bekannt sind, aber nur begrenzt als alltägliches Zahlungsmittel genutzt werden. Viele Nutzer halten sie eher als Anlage oder Spekulationsobjekt. Dabei ist zu beachten, dass die Daten auf Selbstauskünften beruhen und nicht jede einzelne Transaktion messen~\cite{Zamora2026}.

Der EU-Fall verdeutlicht damit drei Grenzen: Der Staat behält die Kontrolle über das Geldsystem, Regulierung macht Krypto-Assets nicht automatisch zu Geld, und das Vertrauen der Nutzer reicht für alltägliche Zahlungen bisher nicht aus. Europa lässt sich daher als Beispiel für die Strategie „Regulierung statt Einführung“ einordnen.

\endinput